% ============================================================
% EvidenceFirst draft
% Springer LNICST / LNCS style
% Compile from paper/: pdflatex -> bibtex -> pdflatex -> pdflatex
% ============================================================
\documentclass[runningheads]{llncs}

\usepackage[T1]{fontenc}
\usepackage{amsmath,amssymb}
\usepackage{booktabs}
\usepackage{graphicx}
\usepackage{multirow}
\usepackage{xcolor}
\usepackage{xspace}
\usepackage[hidelinks]{hyperref}
\usepackage{microtype}
\usepackage{algorithm}
\usepackage{algpseudocode}

\newcommand{\sys}{\textsc{EvidenceFirst}\xspace}
\newcommand{\eg}{\emph{e.g.},\xspace}
\newcommand{\ie}{\emph{i.e.},\xspace}
\newcommand{\etal}{\emph{et al.}\xspace}
\newcommand{\calG}{\mathcal{G}}
\newcommand{\calQ}{\mathcal{Q}}
\newcommand{\calE}{\mathcal{E}}
\newcommand{\calC}{\mathcal{C}}
\newcommand{\calP}{\mathcal{P}}

\begin{document}

\title{\sys: Pre-Generation Evidence-Chain Verification for Traceable KG-RAG}
\titlerunning{\sys: Pre-Generation Evidence-Chain Verification}

\author{Anonymous Author(s)}
\authorrunning{Anonymous}
\institute{
  Institution anonymized for review\\
  \email{email@anonymized.com}
}

\maketitle

% ============================================================
\begin{abstract}
% ============================================================

Knowledge graph-enhanced retrieval-augmented generation (KG-RAG) is attractive
for multi-hop question answering because graph structure can expose the entity
and relation chain needed to answer a question. Yet most KG-RAG systems still
verify too late. They retrieve a graph neighborhood, ask a language model to
generate an answer, and only then evaluate the final string. When the retrieved
subgraph is incomplete, the generator may still produce a fluent answer, but
the system cannot say whether the failure came from a missing entity, a broken
bridge relation, a disconnected component, or the generator itself.

We present \sys, a training-free KG-RAG pipeline that verifies evidence before
generation. The central object of verification is not the answer string, but
the graph-structured evidence chain. \sys builds a per-question evidence graph,
checks whether the graph contains a complete multi-hop chain, diagnoses graph
gaps, performs localized repair, and then generates a grounded answer from the
verified or repaired evidence. The output is therefore both an answer and an
audit trail: chain completeness, missing entities, disconnected pairs, repair
attempts, fallback cases, and runtime failures.

On a repaired-context 1000-question HotpotQA evaluation, the frozen fresh
\sys run reaches 55.0\% EM and 64.8\% F1, competitive with the strongest
standalone HopRAG baseline (54.7\% EM / 65.0\% F1). Deterministic offline
answer cleanup and selector guarding reaches 56.2\% EM and 65.6\% F1 without
additional LLM calls. We do not frame \sys as a score-only state-of-the-art
claim. The contribution is that pre-generation evidence-chain verification can
match strong standalone baselines while exposing graph-level diagnostics that
retrieve-then-generate systems do not provide.

\keywords{KG-RAG \and GraphRAG \and Evidence Verification \and Multi-Hop QA
\and HotpotQA \and Traceable Reasoning}
\end{abstract}

% ============================================================
\section{Introduction}
\label{sec:intro}
% ============================================================

Multi-hop question answering over Web knowledge requires more than retrieving
a passage that shares surface words with the question. A correct answer often
depends on a chain of intermediate entities and relations: an entity mentioned
in the question leads to a bridge page, the bridge page identifies another
entity, and only then can the answer be derived. Retrieval-augmented generation
(RAG) grounds a language model in retrieved text~\cite{lewis2020rag}; KG-RAG
and GraphRAG systems add structured entity-relation evidence over the retrieved
corpus~\cite{edge2024graphrag,guo2024lightrag,han2025graphragsurvey}. Graphs
are a natural fit for multi-hop reasoning because the answer is often the
endpoint of a path through evidence.

The open problem is not merely retrieval quality. It is evidence sufficiency.
A KG-RAG system may retrieve several relevant triples and still miss the bridge
edge that makes the answer derivable. It may retrieve both endpoint entities
but leave them in disconnected components. It may retrieve a graph neighborhood
whose entities look plausible but do not form the required reasoning chain.
In all of these cases, a retrieve-then-generate pipeline can still produce a
confident answer. The final answer string may be evaluated later, but the
system has already lost the ability to explain whether the failure was caused
by retrieval, graph construction, evidence connectivity, or generation.

\sys addresses this by making the evidence graph the verified object. Before
generation, the system checks whether the retrieved graph contains a complete
multi-hop evidence chain for the question. If the chain is incomplete, \sys
does not simply lower confidence. It emits a graph-structural diagnostic:
which entity is missing, which endpoints are disconnected, whether a bridge
edge is absent, or whether evidence retrieval failed entirely. The diagnostic
then drives targeted repair, after which the generator receives the verified
or repaired evidence.

This framing deliberately changes the paper's goal. We do not argue that every
new RAG paper should be judged only by an incremental EM improvement. Instead,
we ask a stricter systems question: can a KG-RAG pipeline remain competitive
with strong standalone baselines while providing an audit trail for evidence
sufficiency? On the current repaired-context HotpotQA evaluation, the answer is
yes. The frozen fresh \sys run reaches 55.0\% EM and 64.8\% F1, comparable to
the strongest standalone HopRAG result. A deterministic offline cleanup pass
raises this to 56.2\% EM and 65.6\% F1, but the paper keeps this boundary
explicit: the fresh run is the primary standalone result, while offline cleanup
is reported as engineering analysis.

\paragraph{Contributions.}
This paper makes three contributions:
\begin{itemize}
    \item We introduce pre-generation evidence-chain verification for KG-RAG:
    the verified object is the retrieved evidence path, not only the generated
    answer.
    \item We define graph-structural diagnostics for multi-hop QA, including
    chain completeness, missing entities, disconnected pairs, bridge gaps, and
    fallback cases.
    \item We evaluate \sys on a repaired-context 1000-question HotpotQA setup,
    showing competitive answer quality with HopRAG while preserving a clear
    diagnostic trace for each prediction.
\end{itemize}

% ============================================================
\section{Related Work}
\label{sec:related}
% ============================================================

\paragraph{Retrieval-augmented generation.}
RAG grounds language-model generation in retrieved text passages~\cite{lewis2020rag}. Dense retrieval methods such as DPR~\cite{karpukhin2020dpr}
and iterative systems such as IRCoT~\cite{trivedi2022ircot} improve recall by
retrieving evidence across reasoning steps. Active retrieval and self-reflective
generation methods further adapt retrieval to the model's intermediate state~\cite{jiang2023active,asai2023selfrag}. These methods improve access to
evidence, but they usually do not expose graph-structural reasons why a
multi-hop evidence chain is complete or incomplete.

\paragraph{KG-RAG and GraphRAG.}
GraphRAG systems organize retrieved evidence into entities, relations, and
communities~\cite{edge2024graphrag,guo2024lightrag}. Recent work explores
graph traversal and path-aware retrieval for multi-hop QA~\cite{gutierrez2024hipporag,chen2025pathrag,yu2025hoprag}. These systems
show that graph structure can improve multi-hop retrieval, but the graph is
usually treated as a retrieval substrate rather than an object to be verified.
\sys differs by making the evidence graph itself the target of pre-generation
verification.

\paragraph{Verification and auditability.}
Verification-oriented RAG methods check generated claims or filter unreliable
answers~\cite{fairrag2025,huang2025unbiased}. Agent frameworks also expose
tool traces and intermediate decisions~\cite{wu2023autogen,hong2024metagpt}.
However, answer-level verification can happen too late: once the generator has
filled gaps in the evidence, it is difficult to distinguish a wrong answer
caused by generation from one caused by incomplete retrieval. \sys instead
checks evidence-chain sufficiency before generation and records the graph gap
responsible for incomplete reasoning.

% ============================================================
\section{Method}
\label{sec:method}
% ============================================================

\subsection{Design Principles}
\label{sec:principles}

\sys is built around three design principles.

\paragraph{Verify before generation.}
The system should not ask a generator to answer from evidence whose multi-hop
structure has not been checked. Verification therefore happens before final
answer generation.

\paragraph{Verify the evidence chain, not only the answer.}
Answer-level verification asks whether a candidate answer is plausible or
supported. \sys asks a prior question: whether the retrieved graph contains a
connected evidence chain from question entities to answer-supporting evidence.
This makes the verifier independent of any particular generated answer string.

\paragraph{Expose graph gaps as first-class outputs.}
The verifier returns structured diagnostics, not just a confidence score. A
wrong answer and an incomplete evidence graph are different failure modes; the
system should expose which one occurred.

\subsection{Problem Setup}
\label{sec:setup}

Let $q$ be a multi-hop question and $\calP=\{p_i\}$ a set of candidate
passages. A KG construction step extracts triples
$\calG=(V,E)$, where vertices are entities and edges are typed relations or
text-derived relational statements. The goal is to produce an answer $a$ while
also determining whether the evidence graph contains a structurally complete
chain for $q$.

Unlike answer-verification methods, \sys does not begin by asking whether a
generated answer is correct. It asks whether the evidence graph contains enough
connected support to justify any answer. This changes the key intermediate
state from an answer candidate to an evidence-chain state:
\[
    s(q,\calG) = (\calC, d, r),
\]
where $\calC$ is the retrieved evidence chain, $d$ is a diagnostic label, and
$r$ is a repair action or fallback decision.

\subsection{Pipeline Overview}
\label{sec:pipeline}

\sys uses six stages:
\begin{enumerate}
    \item Build or load a per-question KG from the candidate passages.
    \item Retrieve KG triples and top-ranked passages for the question.
    \item Construct an evidence graph from retrieved triples.
    \item Check graph-structural chain completeness.
    \item Repair localized graph gaps when possible.
    \item Generate a concise answer from the verified or repaired evidence.
\end{enumerate}

Algorithm~\ref{alg:evidencefirst} summarizes the procedure.

\begin{algorithm}[t]
\caption{\sys inference}
\label{alg:evidencefirst}
\begin{algorithmic}[1]
\Require Question $q$, passages $\calP$, per-question KG $\calG$
\Ensure Answer $a$, diagnostic record $D$
\State $T \gets$ retrieve\_triples$(q,\calG)$
\State $G_E \gets$ build\_evidence\_graph$(T)$
\State $(complete, D) \gets$ check\_chain\_completeness$(q,G_E)$
\If{not $complete$}
    \State $(T', D_r) \gets$ repair\_gap$(q,\calG,D)$
    \State $T \gets T \cup T'$
    \State $G_E \gets$ build\_evidence\_graph$(T)$
    \State $(complete, D') \gets$ check\_chain\_completeness$(q,G_E)$
    \State $D \gets D \cup D_r \cup D'$
\EndIf
\If{$T$ is sparse}
    \State $P_q \gets$ BM25\_passage\_fallback$(q,\calP)$
\Else
    \State $P_q \gets \emptyset$
\EndIf
\State $a \gets$ grounded\_generate$(q,T,P_q,D)$
\State $a \gets$ answer\_canonicalize$(q,a,\calP)$
\State \Return $(a,D)$
\end{algorithmic}
\end{algorithm}

\subsection{Evidence-Chain Verification}
\label{sec:verification}

The chain verifier computes whether the evidence graph contains the entities
and connections required by the question. It identifies four gap types.

\begin{table}[t]
\centering
\caption{Graph-structural diagnostics emitted by \sys before generation.}
\label{tab:gap_types}
\begin{tabular}{lll}
\toprule
Gap type & Graph condition & Typical repair action \\
\midrule
Missing entity & Required node absent & Entity expansion / passage recall \\
Disconnected pair & Endpoints in separate components & Bridge search \\
Bridge gap & Partial path without bridge edge & Targeted relation retrieval \\
Empty evidence & No usable triples retrieved & Passage fallback \\
\bottomrule
\end{tabular}
\end{table}

\paragraph{Missing entity.}
A question entity, candidate answer entity, or expected bridge entity is absent
from the retrieved evidence graph. This usually indicates failed entity linking,
insufficient KG extraction, or passage truncation.

\paragraph{Disconnected pair.}
Two required entities are present but lie in disconnected components. This
means retrieval found relevant endpoints but missed the bridge relation.

\paragraph{Bridge gap.}
A partial path exists, but a specific edge or bridge entity is missing. This is
the most directly repairable case, because the system can target retrieval at
the missing connection.

\paragraph{Empty evidence.}
No useful triples are retrieved. In this case \sys falls back to passage-level
evidence rather than hallucinating from an empty graph.

The verifier emits a structured diagnostic record rather than a single scalar
confidence score. The diagnostic record is used both for repair and for later
analysis.

\paragraph{Completeness criterion.}
Operationally, a chain is treated as complete when the retrieved evidence graph
contains the required question entities and a connected support structure that
links them through retrieved triples or repaired evidence. The checker is not
an oracle for semantic truth: it is a structural sufficiency test. This
distinction matters. A complete evidence chain can still lead to a wrong answer
if the generator misreads it, while an incomplete chain is a warning that the
system should repair or fallback before generation.

\subsection{Targeted Gap Repair}
\label{sec:repair}

When the chain is incomplete, \sys does not restart the full pipeline. It uses
the diagnostic label to target repair. Missing-entity cases trigger entity
expansion and passage-backed retrieval; disconnected-pair cases trigger bridge
search; sparse-chain cases trigger passage fallback. This keeps repair scoped
to the graph failure observed by the verifier.

This repair policy is deliberately conservative. A repair attempt can add
evidence, but it does not override the need for grounded generation. If the
graph remains incomplete, the final diagnostic record preserves that fact.

\subsection{Grounded Generation and Answer Canonicalization}
\label{sec:generation}

The answer generator receives the verified chain, repaired triples, and
fallback passages when needed. The prompt asks for a concise final answer and
records whether the answer came from a complete chain, a repaired chain, or a
fallback path.

After generation, \sys applies deterministic answer canonicalization. This is
not used to change the evidence graph; it only normalizes final answer strings
for exact-match evaluation. Examples include removing instruction leakage,
extracting acronym answers from definition-like outputs, keeping supported
full names from passages, and selecting a short final answer segment when the
model emits leaked context before the answer. We report fresh results and
posthoc deterministic cleanup results separately.

% ============================================================
\section{Experimental Setup}
\label{sec:setup_exp}
% ============================================================

\subsection{Dataset}

We evaluate on HotpotQA distractor-style multi-hop questions~\cite{yang2018hotpotqa}. During development we found that the local sample
file used by earlier experiments contained incomplete context: in the first
100 rows, 110 of 240 supporting-fact titles were absent from the provided
context. We therefore restore contexts by question id from the HuggingFace
HotpotQA distractor validation split and use the repaired-context file as the
primary EvidenceFirst evaluation input. On the current 1000-question evaluation
file, all supporting-fact titles are present in context.

The repository also contains a Wikidata-derived 2WikiMultiHopQA sample
~\cite{ho2020constructing}. This 500-question sample stores natural-language
contexts together with gold evidence triples, Wikidata entity ids, evidence
ids, and answer ids. For example, a compositional question can include a gold
chain such as $(\textit{film}, \textit{director}, \textit{person})$ followed by
$(\textit{person}, \textit{place of birth}, \textit{city})$. This makes
2Wiki/Wikidata a particularly good fit for evidence-chain verification because
the graph structure is explicitly available. In the current draft, HotpotQA is
the primary EvidenceFirst result; the 2Wiki/Wikidata sample is reported as an
available secondary benchmark and a natural next rerun target.

\subsection{Baselines}

We compare against the repository's 1000-question HotpotQA baseline summary:
Naive RAG, IRCoT, A-RAG, LightRAG, Microsoft GraphRAG, HopRAG strict rescoring,
and HopRAG legacy metrics. We also report internal selector and consensus
variants for context, but we do not treat external-consensus variants as
standalone systems because they use outputs from other baselines.

For 2Wiki/Wikidata, we report the surviving 500-question legacy CoMaGRAG
baselines only as scope context. These rows are useful because they establish
that the dataset, KG cache, and evaluation pipeline already exist in the
repository. They should not be read as EvidenceFirst results until the
EvidenceFirst pipeline is rerun on the 2Wiki/Wikidata sample.

\subsection{Metrics}

We use standard HotpotQA exact match (EM) and token F1, with normalization by
lowercasing, removing articles, removing punctuation, and collapsing
whitespace. In addition, \sys records evidence diagnostics: chain completion,
repair attempts, missing entities, disconnected pairs, fallback frequency, and
runtime errors.

\subsection{Implementation Notes}

The current run uses Qwen-compatible LLM calls through the repository's OpenAI
interface and a local SentenceTransformer embedding model for graph retrieval.
The 1000-question run was executed in split batches to avoid GPU memory
contention from the local embedding model. The first complete merged result
with no CUDA OOM failures is frozen as v4. Later v5/v6 rows are deterministic
offline passes over the v4 predictions and make no fresh LLM calls.

\paragraph{Reporting boundary.}
We report three kinds of numbers and keep them separate. The v4 row is the
fresh end-to-end \sys run and is the primary method result. The v5/v6 rows are
deterministic posthoc answer-normalization and selector-guard passes over the
same v4 predictions; they measure how much of the remaining error is caused by
formatting or conservative answer selection. The external consensus row is a
hybrid upper comparison that uses agreement among external baseline outputs;
it is not a standalone \sys result.

% ============================================================
\section{Results}
\label{sec:results}
% ============================================================

\subsection{Main HotpotQA Results}

Table~\ref{tab:main_results} compares \sys with the available 1000-question
baselines. The fresh \sys run is competitive with the strongest standalone
graph baseline: it slightly exceeds HopRAG legacy EM and is within 0.2 F1
points. The deterministic offline selector-guard pass improves both EM and F1,
but we report it separately because it is posthoc over the same fresh
predictions.

\begin{table}[t]
\centering
\caption{HotpotQA 1000-question results. v4 is the frozen fresh \sys run.
v6 is deterministic offline cleanup and selector guarding over v4/v5
predictions with no fresh LLM calls.}
\label{tab:main_results}
\begin{tabular}{lccc}
\toprule
Method & N & EM & F1 \\
\midrule
Microsoft GraphRAG & 1000 & 0.1890 & 0.2774 \\
Naive RAG & 1000 & 0.3920 & 0.4865 \\
LightRAG & 1000 & 0.4380 & 0.5376 \\
IRCoT & 1000 & 0.4660 & 0.5707 \\
A-RAG & 1000 & 0.5170 & 0.6273 \\
HopRAG strict & 1000 & 0.5410 & 0.6438 \\
HopRAG legacy & 1000 & 0.5470 & 0.6498 \\
\midrule
\sys v4 fresh & 1000 & 0.5500 & 0.6484 \\
\sys v6 offline selector guard & 1000 & 0.5620 & 0.6564 \\
\midrule
External consensus hybrid & 1000 & 0.5840 & 0.6869 \\
\bottomrule
\end{tabular}
\end{table}

The result supports a conservative claim: \sys is not presented as a
state-of-the-art score-only system, but as a traceable KG-RAG pipeline whose
fresh accuracy is competitive with the strongest standalone baselines. The
offline pass shows that a small number of deterministic answer-selection
failures can be corrected without additional model calls.

Two observations follow from Table~\ref{tab:main_results}. First, the fresh
\sys row is already in the same accuracy regime as HopRAG. This is important
because \sys adds diagnostics without giving up answer quality. Second, the
offline v6 row exceeds the standalone HopRAG rows, but it should be interpreted
as deterministic cleanup over a frozen run, not as a new retrieval or reasoning
experiment. This is why the paper's main claim is competitive accuracy plus
auditability rather than an unrestricted leaderboard claim.

\subsection{Repaired-Context 100-Question Check}

On the repaired-context first 100 HotpotQA rows, \sys reaches 0.5800 EM and
0.6495 F1. This slice is useful as a development check because it is small
enough for fast iteration while preserving complete supporting contexts. The
larger 1000-question v4 result is the primary fresh-run result.

\subsection{Offline Answer Cleanup}

Table~\ref{tab:offline_passes} summarizes the deterministic posthoc passes.
The v5 pass fixes final-answer leakage such as ``context sentence followed by
answer''. The v6 pass adds a narrow selector guard for cases where the current
answer is clearly the wrong type, such as a non-yes/no question answered with
``no'', while the postprocessor produced a short supported phrase.

The no-loss property in Table~\ref{tab:offline_passes} is measured against the
frozen v4/v5 predictions. It does not imply that the rule is universally safe
under distribution shift; rather, it justifies reporting the pass as a
deterministic cleanup of this evaluation artifact.

\begin{table}[t]
\centering
\caption{Deterministic offline passes over the frozen v4 predictions. These
passes do not make fresh LLM calls and are not counted as fresh end-to-end
runs.}
\label{tab:offline_passes}
\begin{tabular}{lcccc}
\toprule
Variant & EM & F1 & EM gains & EM losses \\
\midrule
v4 fresh & 0.5500 & 0.6484 & -- & -- \\
v5 offline final-answer clean & 0.5570 & 0.6507 & 7 & 0 \\
v6 offline selector guard & 0.5620 & 0.6564 & 5 & 0 \\
\bottomrule
\end{tabular}
\end{table}

\subsection{Evidence Diagnostics}

The diagnostic trace is the main reason to use \sys. Table~\ref{tab:diag}
shows the aggregate diagnostics from the frozen v4 1000-question run. More
than half of the questions have complete chains after retrieval and repair,
while a large number still require repair attempts. The disconnected-pair
count shows that many errors are not simply missing isolated entities; they
are failures to retrieve or construct the bridge relation connecting otherwise
present evidence.

\begin{table}[t]
\centering
\caption{EvidenceFirst diagnostics on the frozen v4 1000-question run.}
\label{tab:diag}
\begin{tabular}{lr}
\toprule
Diagnostic & Count \\
\midrule
Questions & 1000 \\
Chain complete & 568 \\
Repair attempts & 637 \\
Fallback cases & 5 \\
Missing-entity total & 432 \\
Disconnected-pair total & 2082 \\
Runtime/content-filter errors & 4 \\
CUDA OOM errors & 0 \\
\bottomrule
\end{tabular}
\end{table}

These diagnostics give a more actionable error surface than answer accuracy
alone. A plain EM failure says only that the final string did not match the
gold answer. A \sys failure can additionally say whether the graph lacked a
question entity, contained disconnected components, or required fallback
because triple retrieval was sparse.

\subsection{2Wiki/Wikidata Scope}

Table~\ref{tab:2wiki_scope} records the existing 500-question
2Wiki/Wikidata results available in the repository. These are legacy CoMaGRAG
rows, not EvidenceFirst rows. We include them in the draft because the dataset
is important for the paper's planned generalization story: unlike HotpotQA, the
2Wiki sample includes gold evidence triples and Wikidata ids, which can support
direct evaluation of evidence-chain completeness.

\begin{table}[t]
\centering
\caption{Existing legacy 2Wiki/Wikidata 500-question results. These rows are
included as dataset scope context, not as EvidenceFirst results.}
\label{tab:2wiki_scope}
\begin{tabular}{lccc}
\toprule
Method & N & EM & F1 \\
\midrule
Naive RAG & 500 & 0.3020 & 0.3361 \\
CoMaGRAG no decomposition & 500 & 0.3040 & 0.3533 \\
CoMaGRAG full & 500 & 0.3180 & 0.3770 \\
CoMaGRAG no verification & 500 & 0.3360 & 0.3862 \\
\bottomrule
\end{tabular}
\end{table}

The current table should be treated as a placeholder for experimental scope,
not as a final result section. The main reason to keep 2Wiki/Wikidata in the
paper is methodological: it provides gold graph evidence that can test whether
\sys's chain verifier identifies the same supporting relation path as the
dataset annotation.

% ============================================================
\section{Discussion}
\label{sec:discussion}
% ============================================================

\paragraph{Accuracy is competitive, but not the main story.}
The fresh \sys result is close to HopRAG, and deterministic cleanup exceeds the
standalone HopRAG numbers available in the repository. However, the strongest
hybrid consensus result remains higher. We therefore avoid framing the method
as a pure leaderboard contribution. The more robust claim is that \sys reaches
competitive accuracy while adding pre-generation evidence diagnostics.

\paragraph{Why pre-generation verification matters.}
Answer-level verification can reject a bad answer, but it often cannot explain
why the system failed. Was the answer wrong because the model reasoned poorly,
because the graph retrieval missed a bridge entity, because entity linking
selected the wrong node, or because the context itself lacked the page? \sys
records the graph state before generation, making these distinctions visible.

\paragraph{Diagnostics turn failures into engineering signals.}
The diagnostic counters are not merely explanatory metadata. They identify
where the next improvement should occur. A high missing-entity rate suggests
better entity extraction or context restoration. A high disconnected-pair rate
suggests bridge retrieval and relation extraction improvements. A high fallback
rate suggests sparse graph construction. This makes \sys useful even when the
final answer is wrong, because the failure is assigned to a repairable system
component.

\paragraph{The context-repair lesson.}
The early local HotpotQA sample had incomplete contexts. This substantially
changed the interpretation of results: a system cannot be expected to recover
supporting pages absent from the evaluation context. Restoring contexts by id
from the original HotpotQA validation split makes the evaluation fairer and
also demonstrates why evidence diagnostics are useful. Missing context and
missing graph evidence should not be conflated.

\paragraph{Why the Wikidata-derived benchmark matters.}
2Wiki/Wikidata is a stronger test of the method's central claim than a plain
answer-only benchmark because its examples include explicit evidence triples
and entity ids. This makes it possible to evaluate not only whether the final
answer string is correct, but whether the system recovers the annotated
relation chain. The current paper draft therefore keeps 2Wiki/Wikidata in view
as the natural next evaluation, while avoiding unsupported EvidenceFirst claims
before the rerun is completed.

\paragraph{Posthoc passes are useful but bounded.}
The v5/v6 deterministic passes improve EM by correcting final-answer formatting
and narrow selector failures. They do not change retrieval or reasoning. In
the paper, these should be presented as analysis and engineering cleanup rather
than as the core method. The core method is the fresh v4 pipeline.

% ============================================================
\section{Limitations}
\label{sec:limitations}
% ============================================================

\sys has four important limitations. First, graph construction quality remains
a bottleneck. If entity extraction or triple extraction fails, the verifier can
diagnose the failure but cannot always repair it. Second, the current repair
policy is conservative and still leaves many disconnected-pair cases unresolved.
Third, v5/v6 improvements are deterministic offline passes over existing
predictions; they should not be reported as fresh LLM runs. Fourth, the current
EvidenceFirst evaluation is focused on HotpotQA. Although the repository
contains a 2Wiki/Wikidata sample and legacy baselines, the EvidenceFirst rerun
on that dataset remains future work. The evidence-chain diagnostics should also
be tested on additional multi-hop datasets and open-domain Web settings.

There are also two threats to validity. The repaired-context evaluation is more
fair than the incomplete local sample, but it depends on restoring contexts by
question id from the source HotpotQA validation split. Future work should
publish the repaired ids and context-generation script alongside the result
artifacts. In addition, some baselines use different retrieval stacks and
runtime assumptions. We therefore treat the comparison as evidence of
competitiveness, not as a definitive benchmark ranking.

% ============================================================
\section{Conclusion}
\label{sec:conclusion}
% ============================================================

We introduced \sys, a KG-RAG pipeline that verifies evidence-chain completeness
before generation. The method changes the reliability object from the final
answer string to the retrieved evidence graph, enabling graph-structural
diagnostics such as missing entities, disconnected pairs, repair attempts, and
fallback cases. On a repaired-context 1000-question HotpotQA evaluation, the
fresh \sys run is competitive with strong standalone baselines, and deterministic
cleanup further improves answer accuracy without new model calls. More
importantly, \sys provides a traceable account of where evidence succeeds or
fails. This makes it a step toward KG-RAG systems that are not only accurate,
but auditable.

\bibliographystyle{splncs04}
\bibliography{comagraag}

\end{document}
